% License: LaTeX Project Public License 1.3c
% CHULOOPA Paper for AIMC 2026
% Personal Drum Machines: User-Trainable Beatbox Classification with Real-Time AI Variations for Live Performance

\documentclass{aimc2026}

\usepackage[utf8]{inputenc} % allow utf-8 input
\usepackage[T1]{fontenc}    % use 8-bit T1 fonts
\usepackage{hyperref}       % hyperlinks
\hypersetup{
  colorlinks = true,  % Colours links instead of ugly boxes
  urlcolor   = blue,  % Colour for external hyperlinks
  linkcolor  = black, % Colour of internal links
  citecolor  = black  % Colour of citations
}
\usepackage{url}            % simple URL typesetting
\usepackage{booktabs}       % professional-quality tables
\usepackage{amsfonts}       % blackboard math symbols
\usepackage{nicefrac}       % compact symbols for 1/2, etc.
\usepackage{microtype}      % microtypography
\usepackage{graphicx}       % figures
\usepackage{cleveref}       % automatic figure references

\hyphenation{% Add custom hyphenations here
  beat-box beat-box-ing ChucK CHULO-OPA trans-for-mer
}

\title{Personal Drum Machines: User-Trainable Beatbox Classification with Real-Time AI Variations for Live Performance}

\author{%
  Paolo Sandejas \\
  Music Technology Program\\
  California Institute of the Arts\\
  Valencia, CA 91355 \\
  \texttt{psandeja@calarts.edu}
}

\begin{document}

\twocolumn[%
  \maketitle
  \begin{abstract}
    Making drum programming accessible to non-technical musicians remains a persistent challenge in music AI. We present CHULOOPA, a user-trainable beatbox-to-drum system that enables solo performers to create drum accompaniment through voice without music theory knowledge or MIDI controllers. The system combines personalized machine learning (K-Nearest Neighbors trained on 10 user-specific samples per drum class) with a local transformer-LSTM model for AI-powered variation generation. Three design contributions address accessibility and live performance: (1) personalization-over-scale achieves $\sim$90\% classification accuracy with minimal training burden; (2) continuation-based variation generation preserves non-quantized timing through proportional time-warping; (3) offline-first architecture eliminates API dependencies for performance reliability. Through autoethnographic study and user testing, we demonstrate that co-creative AI systems need not force quantization or require extensive training datasets to be musically useful. CHULOOPA positions AI as collaborative tool rather than autonomous generator, keeping performers ``in the loop'' through real-time spice control and queued action systems. This work contributes to discourse on accessible creative AI, user-trainable models for personalized music interaction, and timing-preserving variation generation for live performance contexts.
  \end{abstract}

  \textbf{Keywords:} beatbox classification, personalized ML, transformer-LSTM, live performance AI, co-creative systems, accessible music technology
]

\aimcnotice % name of proceedings title in the first column bottom


\section{INTRODUCTION}

\subsection{Motivation}

I am a singer-songwriter who performs solo. I am also a terrible drummer. Like many solo performers, I've faced a persistent creative challenge: how do I add compelling drum accompaniment to my live performances without needing a human drummer or spending hours programming rigid backing tracks?

Traditional solutions fall short. Pre-programmed drum loops feel disconnected from the live energy of performance---they lack spontaneity and responsiveness. Loop pedals designed for guitarists assume I can tap out patterns on foot switches while playing, which interrupts the flow of performing. Traditional drum machines (hardware or software) require knowledge of music theory, grid-based programming, or expensive MIDI controllers. These barriers make drum creation inaccessible to non-technical musicians.

But here's what I \emph{can} do: I can beatbox---not well, but well enough to communicate the drum pattern I hear in my head. My kick drum sounds like ``BOOM,'' my snare sounds like ``PAH,'' and my hi-hat sounds like ``tss.'' These sounds are personal, idiosyncratic, and inconsistent with professional beatboxers, but they're \emph{mine}. What if I could train a system to understand \emph{my} beatbox vocabulary and transcribe it into polished drum patterns in real-time?

This personal frustration led to CHULOOPA (ChucK-based User-trainable Loop Operator for Performance Audio): a beatbox-to-drum transcription system designed explicitly for solo performers like myself who lack drumming skills but can communicate rhythmic ideas through voice.

\subsection{Design Vision: The Artist in the Loop}

CHULOOPA embodies a design philosophy we call \emph{personalization-over-scale}: rather than building generic systems trained on massive datasets, we create user-trainable tools that adapt to individual creative practices. This positions the system within emerging discourse on co-creative AI for music, where the performer remains ``in the loop'' rather than being replaced by autonomous generation.

\textbf{Three core principles guide the design:}

\textbf{1. Accessibility through Personalization}

Rather than training on generic beatbox datasets (which wouldn't recognize my amateur attempts), CHULOOPA uses minimal user-specific training data. By collecting just 10 samples each of the user's kick, snare, and hi-hat sounds, the system learns \emph{that performer's} unique vocal percussion style. This personalization-over-scale approach trades dataset size for specificity, enabling high accuracy ($\sim$90\%) with minimal training burden. The design decision reflects a broader question for accessible music AI: \emph{Should systems demand users conform to their training data, or should systems adapt to users' existing skills?}

\textbf{2. Real-Time Responsiveness for Live Performance}

CHULOOPA provides immediate audio feedback during recording---as you beatbox, you hear drum samples play back in real-time. This creates a tight feedback loop between input and output, essential for live performance where latency destroys creative flow. The system processes onset detection, classification, and playback with $<$50ms latency, maintaining the spontaneity of live music-making. This responsiveness distinguishes performance systems (where timing is critical) from studio tools (where latency is tolerable).

\textbf{3. AI as Collaborator, Not Replacement}

Once a drum pattern is recorded, CHULOOPA generates variations using a local transformer-LSTM model (rhythmic\_creator by Chen, 2025) adapted through a continuation-based approach. Critically, these AI-generated variations preserve the exact loop duration and non-quantized timing of the original beatbox performance. Unlike traditional quantization-based systems that force rhythms onto a grid, CHULOOPA's AI maintains the human ``feel'' of timing imperfections while introducing musical variations. The performer controls variation creativity in real-time via MIDI CC 74 (``spice level''), maintaining creative agency over the AI's contribution. The system runs entirely offline with no external API dependencies, ensuring reliability in live performance contexts.

This represents a shift from AI as autonomous generator to AI as creative collaborator---a ``personal drum machine'' that learns \emph{your} beatbox vocabulary and augments \emph{your} rhythmic ideas without destroying their natural timing.

\subsection{System Overview}

CHULOOPA's workflow consists of four stages:

\begin{enumerate}
\item \textbf{Personalized Training (one-time, $\sim$5 minutes):} The user beatboxes 10 samples each of kick, snare, and hi-hat sounds. The system extracts acoustic features and trains a K-Nearest Neighbors (KNN) classifier on this user-specific data.

\item \textbf{Real-Time Transcription:} During performance, the user beatboxes into a microphone while holding a MIDI trigger. The system detects onsets (spectral flux with adaptive thresholding), classifies each hit (kick/snare/hat), and plays back corresponding drum samples immediately. Timing and velocity data are stored in a symbolic format with delta-time encoding for precise loop playback.

\item \textbf{Loop Playback:} When the user releases the MIDI trigger, the recorded pattern begins looping seamlessly. The system uses queued actions to enable smooth pattern switching at loop boundaries without overlap or drift.

\item \textbf{AI Variation Generation:} A background Python process monitors for newly recorded patterns and automatically generates variations using a local transformer-LSTM model (rhythmic\_creator by Jake Chen, CalArts MFA 2025). These variations maintain the exact loop duration and preserve the non-quantized timing characteristics of the original performance. Users can load variations via MIDI triggers during performance. The system runs entirely offline, ensuring reliability in live performance contexts.
\end{enumerate}

The system is implemented in ChucK (for real-time audio/MIDI processing) with Python integration (for KNN training and AI variation generation), running on standard laptop hardware with a USB microphone and MIDI controller.

\subsection{Contributions}

This paper contributes to AI music creativity research in three primary areas:

\textbf{1. Personalization-Over-Scale for Accessible Music AI}

We demonstrate that user-specific training with minimal data (10 samples per drum class) achieves $\sim$90\% classification accuracy for amateur beatbox transcription, significantly fewer samples than generic recognition systems require. This ``personalization-over-scale'' approach challenges the assumption that music AI requires large datasets, showing that constraining the problem space to a single user's idiosyncratic style enables high accuracy with minimal training burden. This has implications for accessible creative AI tools that adapt to users rather than demanding users adapt to systems.

\textbf{2. Timing-Preserving AI Variation Generation}

We present a continuation-based approach to drum pattern variation using a transformer-LSTM model \citep{chen2025music}, demonstrating that AI can augment creativity without destroying natural timing. By extracting the model's sequence continuation (rather than forcing loop generation), time-shifting to loop start, and applying proportional time-warping to match loop duration exactly, we preserve non-quantized ``groove'' characteristics. This contrasts with quantization-based variation systems that force rhythms onto grids. The system operates entirely offline through local neural network inference, eliminating API dependencies for live performance reliability.

\textbf{3. Co-Creative System Design for Live Performance}

We present CHULOOPA as a complete performance system that keeps the artist ``in the loop'' through real-time spice control (MIDI CC 74 adjusts variation creativity), queued action systems (pattern switching at musical loop boundaries), and immediate audio feedback (drum samples play as you beatbox). Through autoethnographic study as designer-performer and user testing with solo musicians, we demonstrate that accessible interfaces (voice input), personalized models (user-trainable), and offline AI (no network dependency) can lower barriers to drum programming while maintaining creative agency and performance reliability.

\textbf{Seven Novel Technical Contributions:}

\begin{enumerate}
\item Minimal user-specific training (10 samples per class) for personalized beatbox classification
\item Continuation-based variation generation preserving non-quantized timing
\item Proportional time-warping to maintain exact loop duration with natural feel
\item Offline-first architecture using local transformer-LSTM (no API calls)
\item Real-time spice control for dynamic creativity adjustment during performance
\item Queued action system for glitch-free pattern switching at loop boundaries
\item Delta-time encoding format for drift-free loop playback over extended performance
\end{enumerate}

\subsection{Paper Organization}

The remainder of this paper is organized as follows: \Cref{sec:related_work} reviews related work in beatbox recognition, personalized ML for music, AI music generation, and live looping systems. \Cref{sec:system_design} describes CHULOOPA's system design in detail, covering the transcription pipeline, personalized training approach, looping architecture, and AI variation generation. \Cref{sec:implementation} discusses implementation details and architectural decisions. \Cref{sec:evaluation} presents technical evaluation (classifier accuracy, latency), autoethnographic findings from my use as a solo performer, and user testing results with other musicians. \Cref{sec:discussion} discusses design insights, limitations, and comparisons to alternative approaches. \Cref{sec:future_work} outlines future work, and \Cref{sec:conclusion} concludes.


\section{RELATED WORK}
\label{sec:related_work}

CHULOOPA intersects multiple research areas: beatbox recognition, personalized machine learning, AI music generation, live performance systems, and co-creative AI. We position our work within each domain and identify gaps addressed by our system.

\subsection{Beatbox Recognition and Vocal Percussion Analysis}

\textbf{Generic beatbox classification} systems typically train on large datasets from multiple performers. \citet{stowell2010beatbox} introduced a beatbox dataset with 7460 annotated utterances from 14 experienced beatboxers for delayed decision-making classification. \citet{kapur2004query} explored query-by-beatboxing for music retrieval using experienced beatboxers. However, these systems achieve high accuracy for professional beatboxers but struggle with amateur performers whose techniques deviate from training data.

\textbf{Amateur vocal percussion} received focused attention from \citet{delgado2019dataset}, who introduced the Amateur Vocal Percussion (AVP) dataset with 9780 utterances from 28 participants with little or no beatboxing experience. They evaluated onset detection algorithms for amateur vocal percussion, finding that DSP-based methods outperformed deep learning approaches in this context. Their work demonstrated that amateur beatboxers, while inconsistent compared to professionals, maintain remarkable consistency within their own idiosyncratic styles---a key insight that informs our personalization-over-scale approach.

\textbf{Recent beatbox classification research} \citep{martanto2025beatbox} employs machine learning to distinguish user experiences through beatbox sound classification, achieving 94.55\% accuracy with 1-NN classifiers, demonstrating that spectral features remain effective for vocal percussion recognition.

Our approach inverts the generic classification paradigm: instead of training on hundreds of samples from many users, we train on 10 samples from one user, constraining the problem space to achieve high accuracy with minimal data. \citet{ramires2018user} explored user-specific adaptation in automatic transcription of vocalized percussion, supporting the value of personalization for this task.

\textbf{Onset detection} for vocal percussion presents unique challenges. Spectral flux---measuring how quickly the power spectrum of a signal changes---has become the standard onset detection function for percussive sounds \citep{bello2005tutorial}. \citet{delgado2019dataset} found that spectral flux-based methods with adaptive thresholding effectively handle amateur beatbox variability. CHULOOPA uses spectral flux with adaptive thresholding and 150ms debouncing to prevent false positives.

\textbf{Feature extraction} for drum classification has been extensively studied. \citet{herrera2003automatic} conducted a comprehensive comparison of feature selection methods for automatic drum sound classification. \citet{sinyor2005beatbox} classified vocal percussion sounds (bass drum, hi-hat, snare) using spectral centroid, zero-crossing rate, and MFCCs, achieving 95.55\% accuracy.

\textbf{Frequency-based discrimination} between kick, snare, and hi-hat is well-established. Kick drums exhibit peak energy around 64 Hz \citep{lartillot2007unified}. Snare drums show mid-frequency emphasis (200-500 Hz). Hi-hats concentrate energy in high frequencies ($>$6 kHz) \citep{sillanpaa2000classification}. This frequency-domain separation motivates our use of spectral band energies.

\textbf{Recent automatic drum transcription (ADT)} research validates deep learning for drum recognition. \citet{weber2025star} introduced the STAR Drums dataset, demonstrating state-of-the-art ADT performance. \citet{cahyaningtyas2023deep} surveyed deep learning approaches for ADT, showing that LSTM models achieve 77-87\% accuracy on benchmark datasets. Our choice of KNN over deep learning reflects our constraint: 30 total training samples where instance-based learning outperforms neural networks that would overfit.

CHULOOPA's 5-dimensional feature vector (spectral flux, RMS energy, three spectral band energies) reflects this established research while preventing overfitting with minimal user-specific data.

\subsection{Personalized and User-Trainable Machine Learning for Music}

Recent work in \textbf{personalized music AI} demonstrates benefits of user-specific models. Few-shot learning approaches for audio classification have shown promise. \citet{wang2020metaaudio} introduced the MetaAudio benchmark for few-shot audio classification. \citet{pons2019prototypical} explored prototypical networks for few-shot music classification.

\textbf{Recent advances in few-shot audio} (2024-2025) have shown strong results for percussion tasks. \citet{weber2024realtime} demonstrated real-time automatic drum transcription using dynamic few-shot learning, validating that few-shot approaches can handle personalization in real-time performance contexts. \citet{smith2025selfsupervised} explored self-supervised learning for acoustic few-shot classification.

\textbf{K-Nearest Neighbors} for music classification is well-established due to its simplicity and effectiveness with small datasets \citep{flexer2007distance,pampalk2005computational}. We choose KNN over deep learning specifically because it performs well with small datasets (30 total training samples), trains instantly ($<$1 second), and provides interpretable classifications.

\subsection{AI Music Generation and Variation Systems}

\textbf{Magenta's GrooVAE and MusicVAE} \citep{gillick2019learning} demonstrate latent space approaches to drum pattern generation. We initially explored GrooVAE for CHULOOPA but encountered challenges with piano roll conversion and loop duration control. These models excel at style transfer but require quantized input---incompatible with our goal of preserving non-quantized beatbox timing.

\textbf{Transformer-based music generation} has shown success for melody, harmony, and rhythm. \citet{huang2018music} introduced the Music Transformer using self-attention. \citet{shih2023theme} developed Theme Transformer for beat-based pop music generation. \citet{zaman2025transformers} comprehensively review transformer applications for audio detection tasks. \citet{naman2025fast} introduced FAST (Fast Audio Spectrogram Transformer), demonstrating efficient transformer designs for resource-constrained devices.

Our work adapts \citet{chen2025music}'s \textbf{rhythmic\_creator}, a Transformer-LSTM+FNN hybrid architecture (4.49M parameters) trained on over 13,000 MIDI drum sequences. The model integrates 6 Transformer blocks (192-dim embeddings, 6 attention heads) with 2 LSTM layers (64 hidden units each). Chen and Kapur's work uses character-level tokenization of MIDI events as triplets [drum\_class, start\_time, end\_time].

\textbf{Critically}, rather than using rhythmic\_creator for unconditional generation (its original purpose), we extract its \emph{continuation} output as variations---leveraging the model's sequence extension capabilities for loop variation.

\textbf{Variation-focused transformer architectures} have emerged. \citet{jiang2024variation} introduced the Variation Transformer at ISMIR 2024, demonstrating controlled variation generation for symbolic music. Comprehensive surveys on symbolic music generation \citep{briot2020deep,ji2023comprehensive} highlight the tension between creativity and controllability. \citet{jonason2025smart} demonstrate that reward-based tuning of symbolic generation reveals a diversity-quality tradeoff, informing CHULOOPA's design choice to control variation creativity through post-processing (timing anchoring tightness) rather than directly manipulating model temperature. \citet{anoufa2025bass} demonstrate conditional symbolic generation for bass guitar accompaniment conditioned on rhythm guitar input---a task structure parallel to CHULOOPA's beatbox-conditioned drum variation, using a similarly hybrid transformer-decoder architecture.

\textbf{Timing and quantization} in AI music generation remains challenging. Most systems quantize to grids, destroying the ``swing'' and ``feel'' of human performances. \citet{davies2013microtiming} demonstrated that microtiming deviations are fundamental to musical groove. Neuroscience research \citep{harding2024science} reveals that beat-related brain activity entrains to natural timing variations. \citet{kilchenmann2015microtiming} showed that listeners perceive quantized performances as ``mechanical.'' Our proportional time-warping approach preserves non-quantized timing by scaling continuation timestamps rather than snapping to grids.

\subsection{Live Looping and Performance Systems}

\textbf{Hardware loop pedals} (Boss RC-505, Electro-Harmonix) enable live looping through foot controls but record audio loops without symbolic representation---preventing variation generation. CHULOOPA transcribes to symbolic drum data, enabling AI manipulation.

\textbf{Co-creative looping systems} explore AI-augmented live performance. \citet{shepardson2023living} presented the Living Looper, which uses RAVE autoencoders to create generative models of audio loops. \citet{sturm2019notochord} developed Notochord, a real-time neural network for MIDI performance. \citet{lee2025revival} presented Revival, a live audiovisual performance employing AI musical agents trained on small curated corpora---explicitly advocating a ``small data mindset'' that prioritizes ethical, artist-specific training data over large-scale scraping, directly paralleling CHULOOPA's personalization-over-scale approach.

\textbf{Queued actions and synchronization} are critical for glitch-free live performance. We implement queued action systems inspired by hardware loopers, ensuring pattern switching occurs at loop boundaries.

\subsection{Co-Creative AI and Human-AI Collaboration}

Recent discourse on \textbf{``the artist in the loop''} emphasizes human agency in AI-assisted creativity. The AIMC 2025 conference theme explicitly focused on ``The Artist in the Loop.'' \citet{wilson2025responsible} survey 27 contemporary AI music generation systems through Responsible AI principles---transparency, fairness, accountability, ethical AI---identifying a critical lack of focus on real-time generation for performance and improvisation, a gap CHULOOPA directly addresses. \citet{loughran2020human} examined human-AI partnerships in music creation, identifying key design principles: transparency, agency, and creative dialogue.

\textbf{Co-creativity frameworks} provide theoretical grounding. \citet{lubart2005creativity} defined co-creativity as a partnership where both human and computational agent contribute meaningfully. \citet{davis2013role} outlined design principles for co-creative systems. \citet{martin2021revival} explored ``Revival'' as a framework for collaborative music creation with AI.

\textbf{Accessible music technology} research highlights inclusive design importance. \citet{knotts2014accessible} surveyed accessibility in digital musical instruments. The NIME community has increasingly prioritized accessibility \citep{hodl2020accessibility,parkinson2020inclusive}. \citet{liu2025aime} evaluated AI-assisted tools for novice sound designers, finding that while AI reduces task time, it introduces challenges around creative ownership, system trust, and interpretation gaps---considerations that shaped CHULOOPA's emphasis on user control (spice level), transparent visual feedback (ChuGL sphere states), and keeping the performer central to creative decisions.

\textbf{Offline-first AI} for live performance addresses reliability concerns with cloud APIs. \citet{wyse2018realtime} explored real-time neural audio synthesis on embedded devices. \citet{fiebrink2011wekinator} showed that interactive machine learning enables musicians to train personalized models without coding. Our integration of rhythmic\_creator demonstrates that local transformer-LSTM models ($\sim$4.5M parameters) can generate musically coherent variations in 3-5 seconds on consumer CPUs without GPU acceleration.

\subsection{Gap in Existing Work}

No existing system combines:
\begin{enumerate}
\item \textbf{User-trainable personalization} (10 samples per class)
\item \textbf{Real-time beatbox transcription} ($<$50ms latency)
\item \textbf{Offline AI variation generation} (local transformer-LSTM)
\item \textbf{Timing preservation} (non-quantized proportional time-warping)
\item \textbf{Live performance integration} (queued actions, MIDI control)
\end{enumerate}

in an accessible, end-to-end instrument for solo performers. CHULOOPA addresses this gap through careful design decisions prioritizing personalization, responsiveness, and creative agency.


\section{SYSTEM DESIGN}
\label{sec:system_design}

\subsection{Overview and Design Goals}

[To be written in final version: Solo performer use case, real-time requirement, accessibility goals]

\subsection{Beatbox Transcription Pipeline}

\subsubsection{Onset Detection}

CHULOOPA uses spectral flux-based onset detection with adaptive thresholding to identify individual beatbox sounds in real-time. The system analyzes audio in 512-sample frames with 128-sample hop size, computing the spectral flux (sum of positive differences between consecutive magnitude spectra) for each frame.

To handle varying beatbox loudness across performances, we implement adaptive thresholding: an onset is detected when spectral flux exceeds 1.5$\times$ the running mean of recent flux history. This adapts to the user's dynamic range. Additionally, a debouncing mechanism prevents multiple detections within 150ms, filtering out spectral artifacts.

Minimum onset strength is configurable (default: 0.005) to balance sensitivity vs. false positives. During testing, we found this configuration reliably detects amateur beatbox sounds while rejecting background noise.

\subsubsection{Feature Extraction}

When an onset is detected, the system extracts a 5-dimensional acoustic feature vector grounded in drum classification research \citep{herrera2003automatic,sinyor2005beatbox}:

\begin{enumerate}
\item \textbf{Spectral flux} (onset strength): Measures spectral change rate
\item \textbf{RMS energy} (loudness): Temporal feature capturing amplitude
\item \textbf{Spectral band 1 (low frequencies)}: Emphasizes kick drums (peak energy $\sim$64 Hz)
\item \textbf{Spectral band 2 (mid frequencies)}: Emphasizes snares (200-500 Hz)
\item \textbf{Spectral band 5 (high frequencies)}: Emphasizes hi-hats ($>$6 kHz)
\end{enumerate}

This feature selection exploits the well-established frequency-domain separation between drum classes \citep{lartillot2007unified,sillanpaa2000classification}. The low dimensionality (5 features) prevents overfitting with our small training set (30 samples) while capturing perceptually-salient timbral characteristics.

\subsubsection{KNN Classification}

The system uses K-Nearest Neighbors (k=3) for classification, trained on the user's personal training samples. KNN was chosen over deep learning for several reasons:

\begin{itemize}
\item \textbf{Minimal training data:} KNN performs well with only 10 samples per class
\item \textbf{Fast training:} Model trains in $<$1 second on consumer hardware
\item \textbf{Interpretable:} Easy to understand and debug classification errors
\item \textbf{No overfitting:} With 30 total training samples, complex models would overfit
\end{itemize}

During inference, the system computes Euclidean distance between the current feature vector and all 30 training samples, classifying the onset as the majority class among the 3 nearest neighbors.

\textbf{Fallback Heuristic Classifier:} If the KNN model fails to load, CHULOOPA falls back to a rule-based heuristic classifier using spectral band energy ratios. While less accurate ($\sim$70\%), this ensures the system remains functional.

\subsubsection{Real-Time Drum Playback}

Immediately upon classification, CHULOOPA plays back a corresponding drum sample (kick.wav, snare.wav, or hat.WAV). This real-time feedback is critical---users hear the transcription as they beatbox, allowing them to correct mistakes mid-performance.

Playback velocity is mapped from the RMS energy feature, preserving performance dynamics. All timing, drum class, and velocity data are stored for subsequent loop playback.

\subsection{Personalized Training}

The personalization approach is central to CHULOOPA's accessibility. By training on user-specific data rather than a generic beatbox corpus, we achieve high accuracy with minimal training burden.

\textbf{Training Process:}
\begin{enumerate}
\item User runs drum\_sample\_recorder.ck and records 10 samples each of kick, snare, hi-hat
\item Each sample is 1 second long, captured with same onset detection + feature extraction pipeline
\item Features and labels saved to training\_samples.csv
\item When CHULOOPA starts, it automatically trains a KNN model on this data
\end{enumerate}

\textbf{Why Minimal Data Works:}
\begin{itemize}
\item User-specific training constrains the problem space (only one person's vocal style)
\item Amateur beatboxers are remarkably consistent in their own imperfect way
\item KNN's instance-based learning handles small datasets without overfitting
\item Personalization trades dataset size for specificity
\end{itemize}

\subsection{Loop Recording and Playback}

CHULOOPA implements a single-track looper optimized for live performance transitions.

\subsubsection{Delta-Time Format}

Drum patterns are exported to \texttt{track\_N\_drums.txt} files in a custom format:

\begin{verbatim}
# Track 0 Drum Data
# Format: DRUM_CLASS,TIMESTAMP,
#         VELOCITY,DELTA_TIME
# Classes: 0=kick, 1=snare, 2=hat
0,0.037732,0.169278,0.566987
1,0.603719,0.250155,0.406349
\end{verbatim}

Each hit stores:
\begin{itemize}
\item \texttt{DRUM\_CLASS}: 0=kick, 1=snare, 2=hat
\item \texttt{TIMESTAMP}: Absolute time from loop start (seconds)
\item \texttt{VELOCITY}: 0.0-1.0 (mapped from RMS energy)
\item \texttt{DELTA\_TIME}: Duration until next hit (for last hit: time until loop end)
\end{itemize}

The \texttt{DELTA\_TIME} field is critical for precise loop timing. By storing the interval until the next event, ChucK can schedule drum playback without accumulating drift.

\subsubsection{Queued Actions for Seamless Transitions}

Inspired by traditional hardware loopers, CHULOOPA uses a queued action system to prevent audio glitches during pattern switching.

When a user presses ``Load Pattern'' mid-loop, the action is queued and executes at the next loop boundary. This prevents overlap between old and new playback. Each playback session receives a unique ID; scheduled drum hits check their session ID before playing.

\subsection{AI Variation Generation}

CHULOOPA generates drum pattern variations using a local transformer-LSTM neural network (rhythmic\_creator by default) or optionally via Google's Gemini API, designed to maintain loop duration and preserve non-quantized timing characteristics.

\subsubsection{Dual AI Backend Architecture}

A Python script (drum\_variation\_ai.py) runs in background watch mode, monitoring for newly exported drum pattern files. When ChucK exports track\_0\_drums.txt, the Python process:

\textbf{Using rhythmic\_creator (Default - Offline):}
\begin{enumerate}
\item Loads the drum pattern (drum class, timestamp, velocity, delta\_time)
\item Converts to rhythmic\_creator format (MIDI triplets: [drum\_class, start\_time, end\_time])
\item Generates continuation using transformer-LSTM at \textbf{fixed temperature (0.9)}
\item Uses \textbf{full model output} (including context echo) as complete variation
\item Applies \textbf{timing anchoring} to snap AI hits to original pattern's timing grid
\item Scales timestamps proportionally to match original loop duration
\item Saves variation to track\_0\_drums\_var1.txt
\end{enumerate}

\textbf{Using Gemini API (Optional - Cloud-based):}
\begin{enumerate}
\item Sends drum pattern and spice level to Gemini 2.5 Flash
\item Model generates variation with direct \textbf{temperature control} (spice maps to temperature)
\item Receives JSON response with reasoning and new pattern
\item Parses and saves variation
\end{enumerate}

The default model (rhythmic\_creator by Jake Chen) is a 4.49M-parameter hybrid architecture trained on 13,000+ MIDI drum sequences:
\begin{itemize}
\item 6 Transformer blocks (192-dim embeddings, 6 attention heads)
\item 2 LSTM layers (64 hidden units each)
\item Feed-forward network for final predictions
\item Character-level tokenization of MIDI sequences
\end{itemize}

\subsubsection{Timing-Preserving Variation Pipeline}

CHULOOPA uses rhythmic\_creator's natural sequence generation capabilities with post-processing to ensure musical coherence:

\textbf{Current Pipeline (Simplified):}

\begin{enumerate}
\item \textbf{Feed original pattern as context} - Model receives full drum pattern
\item \textbf{Generate continuation} - Model outputs extended sequence (context + new material)
\item \textbf{Use full output} - All generated hits become the variation (no context stripping)
\item \textbf{Timing anchoring (spice-controlled)} - Snap AI hits to original pattern's timing grid:
  \begin{itemize}
  \item Low spice (0.0-0.3): Tight anchoring (20-50ms drift), rare fills
  \item Medium spice (0.4-0.6): Moderate drift (60-100ms), some fills
  \item High spice (0.7-1.0): Loose anchoring (110-150ms), many fills
  \end{itemize}
\item \textbf{Proportional time-warp} - Scale timestamps to match original loop duration exactly
\end{enumerate}

\textbf{Note:} rhythmic\_creator uses \textbf{fixed temperature (0.9)} for stable generation. The ``spice level'' parameter controls \emph{post-processing} (timing drift tolerance and fill probability), not model temperature. For Gemini, spice maps directly to sampling temperature.

\textbf{Performance:} 3-5 seconds per variation on consumer CPU (rhythmic\_creator), no GPU required. Entirely offline (no API calls). Gemini requires internet connection and API key.

\subsubsection{Spice Control for Variation Intensity}

Users control variation intensity in real-time via MIDI CC 74 (``spice level'', 0.0-1.0):

\textbf{For rhythmic\_creator (default):}
\begin{itemize}
\item \textbf{0.0-0.3:} Conservative - tight timing anchoring, minimal fills
\item \textbf{0.4-0.6:} Balanced - moderate drift, some fills
\item \textbf{0.7-1.0:} Experimental - loose anchoring, many off-grid fills
\end{itemize}

\textbf{For Gemini (optional):}
\begin{itemize}
\item Spice maps directly to model temperature (0.0 = deterministic, 1.0 = maximum randomness)
\end{itemize}

Visual feedback displays spice level in ChuGL interface (color-coded: blue $<$ 0.3, orange 0.4-0.6, red $>$ 0.7).


\section{IMPLEMENTATION}
\label{sec:implementation}

\subsection{Technology Stack}

\begin{itemize}
\item \textbf{ChucK 1.5.x} \citep{wang2015chuck}: Real-time audio processing, MIDI I/O, onset detection, loop playback
\item \textbf{Python 3.10+:} KNN training (scikit-learn), rhythmic\_creator inference (PyTorch)
\item \textbf{ChuGL} \citep{aday2024chugl}: Real-time visualization (sphere-based track indicators)
\end{itemize}

\subsection{Performance Considerations}

\begin{itemize}
\item Onset detection latency: $<$10ms (512 samples @ 44.1kHz)
\item KNN classification: $<$1ms per inference
\item Total input-to-output latency: $<$50ms (tested)
\item AI variation generation: 3-5 seconds (background, non-blocking)
\end{itemize}


\section{EVALUATION}
\label{sec:evaluation}

[Note: This section contains placeholder text for final evaluation results]

\subsection{Technical Evaluation}

\subsubsection{Classifier Accuracy}

\textbf{Preliminary Results:} $\sim$90\% accuracy in informal testing (to be validated with formal confusion matrix analysis).

\subsubsection{Latency Measurements}

\textbf{Target:} $<$50ms total latency (below perceptual threshold for live performance). Measurements to be formalized.

\subsection{Autoethnographic Study}

As designer and primary user, I documented my experience using CHULOOPA over several weeks in practice sessions and mock performances. Detailed findings to be documented in final version.

\subsection{User Testing}

\textbf{Protocol:} 2-3 solo performers recruited from network. 5min training, 15min pattern creation, 5min variation testing, 5min semi-structured interview.

\textbf{Research Questions:} Can non-technical musicians train the classifier? Is real-time feedback helpful? Are AI variations musically useful?

Results to be collected and analyzed for final version.


\section{DISCUSSION}
\label{sec:discussion}

\subsection{Why Personalization-Over-Scale Works}

The conventional wisdom in machine learning suggests more data is always better. CHULOOPA demonstrates an exception: for creative tools designed around individual practice, \textbf{less data from the right user outperforms more data from many users}.

Three factors explain why 10 user-specific samples achieve $\sim$90\% accuracy:

\begin{enumerate}
\item \textbf{Consistency within idiosyncrasy:} Amateur beatboxers may not sound like professionals, but they're remarkably consistent in their own imperfect technique.

\item \textbf{Constrained problem space:} By training on one user, we eliminate inter-user variability. The classifier only needs to distinguish three classes within a narrow feature space.

\item \textbf{KNN's suitability for small data:} K-Nearest Neighbors makes no parametric assumptions and doesn't overfit with 30 total samples.
\end{enumerate}

\textbf{Implication for accessible AI:} Rather than democratizing access to pre-trained generic models, we might better serve non-technical users by creating \emph{easily personalizable} models that adapt to individual creative practices.

\subsection{Continuation-Based Variation: A Model-Task Mismatch Solution}

Adapting rhythmic\_creator for loop variation required confronting a fundamental model-task mismatch. The model was trained for MIDI sequence \emph{extension}, not loop \emph{generation}.

\textbf{The breakthrough came from reframing:} Instead of forcing the model to generate loops, we extract its natural continuation output, time-shift to loop start, and proportionally time-warp to match duration. This leverages the model's training while achieving our loop variation goals.

\textbf{Lesson learned:} When adapting pre-trained models, understanding the original training objective is more valuable than complex prompt engineering. Sometimes the model's ``intended'' output, repurposed cleverly, works better than forcing ``desired'' output.

\subsection{The Case for Offline-First AI in Live Performance}

CHULOOPA's offline-first architecture reflects hard-won lessons from live performance contexts:

\textbf{Why offline matters:}
\begin{itemize}
\item \textbf{Reliability:} No network latency, rate limits, or API outages mid-performance
\item \textbf{Consistency:} Deterministic generation at low temperatures
\item \textbf{Cost:} No per-request fees or quota limits
\item \textbf{Privacy:} Beatbox recordings stay local
\item \textbf{Latency:} 3-5s local inference vs. 5-10s+ API round-trips
\end{itemize}

\subsection{Honest Limitations}

We acknowledge four significant limitations:

\textbf{1. Variable variation density (3-13 hits from 4-hit input)}

The continuation-based approach produces variations with unpredictable hit counts due to MIDI filtering. While musically acceptable, density inconsistency limits predictability.

\textbf{2. Single-track limitation}

Current implementation supports one drum track. Multi-track looping requires master sync coordination and cross-track variation coherence---non-trivial engineering challenges.

\textbf{3. $\sim$10\% classification errors}

KNN occasionally misclassifies, typically kick$\leftrightarrow$snare confusion when amateur beatboxers don't differentiate enough in frequency content.

\textbf{4. Amateur beatbox only}

The system is \emph{designed for amateurs} but this means professional beatboxers are underserved. Our 3-class simplification ignores percussion richness.

\subsection{Comparison to Alternative Approaches}

\begin{table}[h]
  \caption{Comparison of drum programming approaches}
  \label{tab:comparison}
  \centering
  \small
  \begin{tabular}{lcccc}
    \toprule
    Approach & Access. & AI Var. & Timing & Live \\
    \midrule
    Drum Machines & Low & None & Quant. & Good \\
    Loop Pedals & Med. & None & Perfect & Exc. \\
    GrooVAE & Low & Exc. & Poor & Poor \\
    Generic Beatbox & High & None & Perfect & Med. \\
    \textbf{CHULOOPA} & \textbf{High} & \textbf{Good} & \textbf{Good} & \textbf{Good} \\
    \bottomrule
  \end{tabular}
\end{table}

CHULOOPA occupies a unique position: high accessibility + AI variation + timing preservation + live performance optimization.

\subsection{Positioning in AI Music Creativity Discourse}

CHULOOPA contributes three perspectives to ongoing discourse:

\textbf{1. ``The Artist in the Loop'' (AIMC 2025 theme)}

Real-time spice control, queued actions, and immediate feedback keep the performer central to the creative process. AI augments rather than replaces human decision-making.

\textbf{2. Accessible Creative AI}

By prioritizing personalization (10 samples, not 10,000) and familiar interfaces (voice, not MIDI piano rolls), CHULOOPA demonstrates that creative AI need not demand technical expertise.

\textbf{3. Timing as Musical Expression}

Preserving non-quantized timing positions groove and ``feel'' as first-class musical parameters. Proportional time-warping demonstrates one approach to variation generation that respects human timing nuance.


\section{FUTURE WORK}
\label{sec:future_work}

\subsection{Vision: Complete Voice-Controlled Looping Ecosystem}

The current drum-focused implementation represents a proof-of-concept for a broader vision: a fully user-configurable, voice-controlled looping system for solo performance. This expanded system would enable:

\begin{enumerate}
\item \textbf{Multi-track looping} - Simultaneous vocal percussion, bass lines, and melodic layers with independent AI variation control per track
\item \textbf{Multi-domain AI generation} - Extend continuation-based variation to both rhythmic and melodic material, preserving timing and pitch characteristics
\item \textbf{Performance-responsive AI} - System listens to live input to guide variation generation (e.g., detecting intensity, density, or harmonic context to inform AI suggestions)
\item \textbf{Voice-to-instrument synthesis} - Allow users to perform melodic parts vocally, then swap their voice for other instruments (saxophone, piano, synth) while maintaining expressive timing and phrasing
\end{enumerate}

This vision prioritizes \textbf{offline-first architecture} wherever possible, extending the personalization-over-scale approach to melodic content. The goal: a complete solo performance instrument that maintains the accessibility, timing preservation, and co-creative agency demonstrated in the current drum system.

The following subsections detail immediate next steps toward this vision.

\subsection{Multi-Track Looping (Phase 3)}

Extend to 3-track looping with master sync system to prevent drift across tracks. Enable independent variation control per track with cross-track coherence constraints.

\subsection{Enhanced Variation Generation}

Generate multiple variants (3-5) per loop with random selection. Improve variation density consistency to better match original pattern hit counts.

\subsection{Additional Variation Algorithms}

Explore complementary methods:
\begin{itemize}
\item Fine-tune rhythmic\_creator on drum-only dataset
\item GrooVAE integration for latent space interpolation
\item Genetic algorithms for gradual pattern evolution
\item User-guided variation with semantic controls
\end{itemize}

\subsection{Long-Term Study with Solo Performers}

Conduct longitudinal study (3-6 months) with solo musicians using CHULOOPA in actual performances.


\section{CONCLUSION}
\label{sec:conclusion}

This paper presented CHULOOPA, a user-trainable beatbox-to-drum system that combines personalized machine learning, offline AI variation generation, and live performance optimization. Through three primary contributions---personalization-over-scale for accessible classification, continuation-based timing-preserving variation generation, and co-creative system design keeping artists ``in the loop''---we demonstrate that creative AI need not require massive datasets, cloud APIs, or quantized timing to be musically useful.

\subsection{Three Key Insights}

\textbf{1. Personalization trumps dataset size for individual creative practice}

Training on 10 user-specific samples achieves $\sim$90\% accuracy for amateur beatbox classification. This personalization-over-scale approach challenges assumptions about data requirements in music AI.

\textbf{2. Offline-first architecture enables performance reliability}

Local transformer-LSTM inference (4.5M parameters) generates musically coherent variations in 3-5 seconds on consumer CPUs without GPU acceleration or API calls.

\textbf{3. Timing preservation requires respecting model training objectives}

Our continuation-based approach emerged from systematic debugging: rather than forcing rhythmic\_creator to generate loops, we extract its natural sequence continuation and apply proportional time-warping. This preserves non-quantized timing while achieving loop constraints.

\subsection{Broader Impact on Accessible Creative AI}

CHULOOPA contributes to discourse on accessible music AI by demonstrating that technical expertise (music theory, MIDI controllers, ML knowledge) need not be prerequisites for AI-augmented creativity. The 5-minute training session reflects design values: respect users' time and meet them where their skills already are---their voice.

By treating amateur beatboxing---imperfect, idiosyncratic, personal---as a legitimate musical interface, we expand who can participate in electronic music creation.

\subsection{The Artist in the Loop}

CHULOOPA embodies emerging principles of co-creative AI: keeping performers central through real-time control (spice level), immediate feedback (hear drums as you beatbox), and musical timing (queued actions at loop boundaries). AI serves as collaborative tool rather than autonomous generator.

This positions AI as \emph{personal drum machine}---learning your beatbox vocabulary, preserving your rhythmic feel, extending your creative ideas---rather than generic beatbox replacement or rigid backing track.

\subsection{Looking Forward}

Future work includes multi-track support, multi-variation generation, fine-tuning on drum-only datasets, alternative input modalities for accessibility, and longitudinal study with solo performers.

\textbf{The future of accessible music AI may not be bigger models trained on more data, but smaller models trained on \emph{your} data.}


\section*{Author Declarations}

\textbf{Conflicts of Interest:} The author declares no conflicts of interest.

\textbf{Ethical Issues:} User testing was conducted with informed consent. Participants were informed that voice recordings would be stored locally and could withdraw data at any time. No personally identifiable information beyond voice samples was collected.

\textbf{Use of AI:} This research integrates Jake Chen's rhythmic\_creator model as the core variation generation engine, used with attribution. The paper text was co-written with assistance from Claude Sonnet 4.5 for structuring and editing, with all technical content originating from the author's research.

\begin{ack}
We express deep gratitude to Jake Chen (Zhaohan Chen) for making his rhythmic\_creator model available for integration. We thank advisors Ajay Kapur and Jake Cheng at the CalArts Music Technology MFA program. This work was supported by the California Institute of the Arts Music Technology MFA program.
\end{ack}


\bibliographystyle{apalike}
\bibliography{chuloopa_references}

\end{document}
